% KHÔNG \input file này: nhiều đoạn bám tài liệu hãng (Hooks, Context, Fetch,
% Docker multi-stage) --- dễ trùng Turnitin. Chương 2 đã đủ chỗ dùng trong đồ án.

\section{Bổ sung chi tiết các công nghệ còn lại}

Các mục dưới đây viết theo cùng khuôn: khái niệm là gì, đặc điểm kỹ thuật liên
quan, chỗ đồ án gắn vào. Mục đích là độc giả ngành CNTT đọc báo cáo không cần
mở tài liệu hãng, đồng thời không nhầm đồ án với giáo trình tổng quát: mỗi
công nghệ kết thúc bằng quyết định hiện thực.

\subsection{LangChain và vị trí của LangGraph}

LangChain cung cấp trừu tượng gọi LLM, ghép prompt, bộ nhớ, công cụ. LangGraph
là lớp đồ thị trạng thái, hợp khi pipeline có nhánh, lặp và cần nhật ký từng
bước. Đồ án dùng LangGraph cho bảy nút hỏi đáp, không dựng agent tự chọn công
cụ tự do (ReAct không giới hạn): pháp lý cần thứ tự truy hồi rồi sinh, không
cần mô hình tự quyết định ``hôm nay có gọi máy tính không''.

Bộ nhớ mặc định \texttt{InMemorySaver} mất khi restart. Đồ án nạp lượt chat từ
PostgreSQL, phù hợp nhiều worker và Docker recreate.

\subsection{Uvicorn}

Uvicorn là máy chủ ASGI. Lệnh phát triển: \texttt{uvicorn} trỏ ứng dụng FastAPI,
reload khi sửa mã. Trong Docker, cùng Uvicorn (không reload) đứng sau cổng 8023.
Không dùng Gunicorn đồng bộ thuần vì SSE. Số worker: đồ án mặc định một tiến
trình trong Compose để Chroma nhúng trong process không bị hai worker đua file
chỉ mục; nhân rộng worker là việc khi tách Chroma thành dịch vụ.

\subsection{asyncpg}

asyncpg là driver PostgreSQL bất đồng bộ. SQLAlchemy 2 dùng nó qua URL
\texttt{postgresql+asyncpg://}. Khác psycopg2 đồng bộ: không chặn vòng sự kiện
khi chờ truy vấn. Chuỗi kết nối, SSL, timeout nằm biến môi trường, không hardcode.

\subsection{Pydantic Settings}

Ngoài schema API, Pydantic đọc biến môi trường vào lớp Settings (prefix, kiểu,
giá trị mặc định). Sai kiểu khóa (ví dụ cổng không phải số) fail sớm lúc khởi
động, tốt hơn fail lúc request đầu. Đồ án tách Settings (bí mật, URL) khỏi
\texttt{config.yaml} (tham số RAG).

\subsection{YAML}

YAML là định dạng cấu hình có lùi hàng, hỗ trợ comment --- khác JSON. Tham số
$k$, trọng số RRF, bật/tắt HyDE nằm YAML để chỉnh không sửa Python. YAML dễ sai
thụt dòng; file đồ án giữ nông, không neo mảng phức tạp.

\subsection{Chroma collection và metadata}

Mỗi Điều khi upsert vào Chroma kèm metadata: \texttt{law\_id}, số Điều, chương,
mã nội bộ. Lọc metadata cho phép hạn chế search space. Kho đồ án chủ yếu một
collection mặc định vì corpus đóng. Xóa collection khi đổi chiều embedding là
bắt buộc: cosine trên hai không gian khác chiều không so được.

\subsection{Cosine, L2 và tích vô hướng}

Khoảng cách L2 phạt độ lớn vector; cosine chuẩn hoá độ dài, phù hợp embedding
đã chuẩn. Chroma cấu hình cosine cho đồ án. Đổi sang L2 mà không đổi cách nhúng
làm lệch thứ tự top-$k$. Công thức cosine đã nêu ở Chương~2; ở đây chỉ nhấn:
metric là một phần hợp đồng chỉ mục, phải ghi vào manifest.

\subsection{Matryoshka representation}

Một số mô hình nhúng huấn luyện để tiền tố vector vẫn mang nghĩa (Matryoshka).
Cắt 1536 từ 3072 chiều dựa trên tính chất này. Không phải mọi mô hình cắt được.
Đổi sang nhà cung cấp khác phải đọc tài liệu mô hình trước khi cắt.

\subsection{Backoff khi gọi API}

Mã 429 (quá tải hạn mức) và 5xx nhà cung cấp xử lý bằng nghỉ dần (exponential
backoff) và giới hạn số lần thử. Đồ án: lô embedding 64 Điều, nghỉ tăng khi 429.
Không thử lại vô hạn: hết lần thì startup/reindex báo lỗi, không treo container
im lặng.

\subsection{React Hooks}

Hook là hàm đặc biệt (\texttt{useState}, \texttt{useEffect}, \texttt{useMemo},
\texttt{useRef}) chỉ gọi ở đỉnh component, cho phép component hàm có trạng thái.
Đồ án: \texttt{useState} cho ô nhập; \texttt{useEffect} khôi phục phiên; hook
TanStack Query bọc \texttt{useQuery}/\texttt{useMutation}. Không dùng class
component. Sai quy tắc hook (gọi trong nhánh \texttt{if}) làm React lệch state
--- các màn hình giữ hook ở top-level.

\subsection{Context API so với Zustand}

Context React truyền dữ liệu sâu cây mà không prop drilling. Context không tối
ưu tần suất cập nhật cao (mọi consumer rerender). Zustand chọn lọc subscribe.
Đồ án: phiên đăng nhập ít trường, đổi thưa --- Zustand đủ. Không bọc toàn app
bằng một Context khổng lồ chứa corpus.

\subsection{Fetch API}

Fetch là API trình duyệt gốc cho HTTP. Đồ án dùng Fetch cho SSE (đọc
\texttt{ReadableStream}), Axios cho REST. EventSource không gắn Bearer. Parser
tách khối \texttt{event:}/\texttt{data:} được Vitest phủ, vì lỗi off-by-one ở
đây làm UI ``đơ'' giữa chừng câu trả lời.

\subsection{FormData và multipart}

Tải PDF dùng \texttt{multipart/form-data}: ranh giới (boundary) tách file và
trường. Axios nếu để \texttt{Content-Type: application/json} thì hỏng. Interceptor
đồ án xóa Content-Type khi body là FormData, để trình duyệt tự gắn boundary.

\subsection{Nginx reverse proxy}

Reverse proxy nhận request từ client, chuyển tiếp server nội bộ, trả đáp. Lợi:
một origin, nén, cache tĩnh, ẩn cổng backend. Cấu hình đồ án: tệp
\texttt{index.html} cho SPA fallback (React Router), \texttt{/api} tới
\texttt{backend:8023}, \texttt{proxy\_buffering off} liên quan SSE --- nếu nginx
đệm hết luồng, UI không thấy token dần.

\subsection{Image Docker đa giai đoạn}

Dockerfile frontend: giai đoạn Node build Vite; giai đoạn nginx copy
\texttt{dist}. Image cuối không chứa npm hay mã nguồn TypeScript, nhỏ hơn và
giảm bề mặt. Backend: image Python 3.10, cài uv, copy mã, không chạy phần tử
dev reload.

\subsection{Volume và mạng cầu Compose}

Volume \texttt{postgres\_data}, \texttt{chroma\_db}, \texttt{uploads} giữ dữ liệu
khi \texttt{docker compose down} (chưa \texttt{-v}). Mạng cầu DNS nội bộ: backend
gọi host \texttt{postgres}, không gọi \texttt{localhost} (localhost trong
container là chính nó). Sai host này là lỗi hay gặp lúc chuyển từ chạy local
sang Compose.

\subsection{Healthcheck}

Compose không chỉ \texttt{depends\_on} khởi động xong process. Healthcheck
Postgres: \texttt{pg\_isready}. Healthcheck backend: HTTP \texttt{/health}.
Backend đợi Postgres healthy mới chạy migration/nạp. \texttt{/health} không đòi
JWT và không gọi Gemini, để vòng lặp Docker không đốt hạn mức.

\subsection{Alembic revision}

Mỗi revision có mã cha. Lịch sử tuyến tính tránh hai nhánh schema. Đồ án không
sửa tay bảng trên container đang chạy rồi quên revision: lần compose máy khác
sẽ thiếu cột. Test tạo database mới chạy upgrade từ đầu.

\subsection{Passlib}

Passlib bọc nhiều thuật toán băm mật khẩu, đồ án chọn argon2. API
\texttt{hash}/\texttt{verify}. Verify sai (hash hỏng, thuật toán lạ) trả False,
không raise --- tránh một hàng CSDL xấu làm cả trang login 500.

\subsection{python-jose hoặc thư viện JWT tương đương}

Tạo và giải mã JWT: header thuật toán, payload claim, chữ ký. Đồ án kiểm
\texttt{exp} và \texttt{type}. Không nhét quyền admin vào localStorage rồi tin
phía client: mọi router admin gọi \texttt{require\_role} trên server.

\subsection{slowapi và khóa đếm}

Đếm theo key user id khi đã đăng nhập, theo IP khi chưa. Chat 20/phút, upload
30/giờ. Số này là tham số, không phải hằng số nghiên cứu. Log 429 giúp phát hiện
vòng lặp frontend (ví dụ poll quá dày --- poll hợp đồng đồ án 2,5 giây, không
phải 50 ms).

\subsection{Vitest và jsdom}

Vitest chạy trên Vite, biến môi trường gần production frontend. jsdom mô phỏng
DOM đủ để test hàm thuần và parser. Không thay thế kiểm thử trình duyệt. 24 test
nhắm chỗ dễ vỡ: SSE, định tuyến splat số hiệu, định dạng ngày task.

\subsection{pytest-asyncio và fixture CSDL}

Test async cần vòng sự kiện. Fixture tạo engine, chạy Alembic, yield, xóa database
tạm. Cấm trỏ fixture vào database demo: một test xóa user sẽ phá dữ liệu người
dùng đang xem.

\subsection{Type hint Python}

PEP 484: chú thích kiểu, công cụ (mypy) kiểm lúc phát triển. FastAPI dùng hint
để sinh OpenAPI. Đồ án không bắt buộc mypy sạch 100\% trên toàn repo trong phạm
vi thời gian, nhưng schema Pydantic --- chỗ biên giới mạng --- được khai báo
bắt buộc.

\subsection{ESLint và ràng buộc TypeScript}

\texttt{strict} TypeScript bắt null. Đồ án hưởng lợi khi field API đổi tên: build
frontend gãy trước khi demo. ESLint không phải phần lõi RAG nhưng giữ import
hook đúng quy tắc.

\subsection{Tailwind JIT}

Just-in-time quét class trong JSX, chỉ phát CSS dùng tới. Class động ghép chuỗi
(\texttt{bg-} + biến) có thể bị bỏ sót; đồ án viết class đủ chữ. Bản 4 đổi chỗ
cấu hình so với bản 3; lock phiên bản trong npm để image build lặp lại được.

\subsection{Khác biệt dev và production frontend}

Dev: Vite HMR, source map, proxy. Production: tệp băm tên (hash) để cache lâu,
không HMR. Cùng một mã gọi \texttt{/api/v1/...} tương đối. Lỗi thường gặp: build
thiếu biến \texttt{VITE\_*} vì biến Vite nhúng lúc build, không lúc chạy nginx.
Đồ án để base URL trống, tránh nhúng \texttt{localhost} vào image.

\subsection{Tách trang quản trị khỏi nghiệp vụ DNNVV}

Router \texttt{/admin} phục vụ quản trị corpus, user và chạy hàng loạt câu hỏi.
Người dùng DNNVV không cần các trang này. Tách route và \texttt{lazy()} giảm
bundle, đồng thời tránh tài khoản thường đốt hạn mức mô hình --- các đường đó
đòi vai trò admin.
